\section{Client Order Flow on Exchange A}
\label{sec:client-flow}

The quoter of Section~\ref{sec:quoter} sets prices, but the only
channel through which those prices become realised revenue is the
stream of client orders reaching Exchange $A$. We therefore need a
generative model of that flow with three layers: \emph{when} an
order arrives, \emph{which type} it is, and \emph{how large} it is.
The intensity baseline $A$, the price sensitivity $k$, the tail
index $\alpha$ of order sizes and the market-order mix $p_{\mathrm{MO}}$
are design parameters of the simulator: they are held fixed across
Phase~1 and Phase~2 and are not re-estimated by the calibration
pipeline of Section~\ref{sec:phase2}. What Phase~2 recalibrates are
the quoter's own risk parameters --- the client-flow knobs above
remain exogenous throughout.

\subsection{Poisson arrivals with price-sensitive intensity}
\label{sec:client-flow-intensity}

Buy-side and sell-side clients are modelled as two independent
inhomogeneous Poisson processes whose instantaneous rates depend on
how far our posted quote sits from the prevailing mid on the
relevant side. Tightening the ask attracts buy-side flow without
mechanically also attracting sell-side flow, which matches the
empirical asymmetry of electronic FX and preserves the separation
of controls the quoter relies on. With the mid $m_t$ and our
top-of-book $(\mathrm{bid}_t, \mathrm{ask}_t)$, the per-side
distances in basis points are
\[
  \delta^{\mathrm{ask}}_t = 10^{4}\,\frac{\mathrm{ask}_t - m_t}{m_t},
  \qquad
  \delta^{\mathrm{bid}}_t = 10^{4}\,\frac{m_t - \mathrm{bid}_t}{m_t},
\]
and the arrival intensities take the exponential kernel of
\citet{avellaneda2008high},
\[
  \lambda^{\mathrm{buy}}_t \;=\; A\,\exp\!\bigl(-k\,\delta^{\mathrm{ask}}_t\bigr),
  \qquad
  \lambda^{\mathrm{sell}}_t \;=\; A\,\exp\!\bigl(-k\,\delta^{\mathrm{bid}}_t\bigr).
\]
Over a simulation step of length $\Delta t$ the number of arrivals
on each side is drawn independently as
$N^{\mathrm{side}} \sim \mathrm{Poisson}(\lambda^{\mathrm{side}}_t \cdot \Delta t)$,
so a step can produce zero, one, or several orders per side. With
the defaults $A = 7\cdot 10^{-3}\,\mathrm{s}^{-1}$ and
$k = 0.3\,\mathrm{bp}^{-1}$, the half-life of the intensity is
$\ln(2)/k \approx 2.3\,\mathrm{bp}$: a quote posted two bps outside
the fair mid already attracts roughly half the arrivals that a
touch-priced quote would.

\begin{figure}[H]
  \centering
  \includegraphics[width=0.75\textwidth]{client-flow-intensity.png}
  \caption{Arrival intensity $\lambda(\delta) = A\,e^{-k\delta}$ as a
           function of the distance from the fair mid, with
           $A = 0.007\,\mathrm{s}^{-1}$ and $k = 0.3\,\mathrm{bp}^{-1}$.
           Each bp of extra distance costs the quote approximately
           $26\%$ of its arrival rate.}
  \label{fig:client-flow-intensity}
\end{figure}

\subsection{Market-versus-limit mix and limit placement}
\label{sec:client-flow-motype}

Each arrival is independently classified market or limit by a
Bernoulli trial with probability $p_{\mathrm{MO}} = 0.5$. Market
orders cross immediately at the touch --- buys at $\mathrm{ask}_t$,
sells at $\mathrm{bid}_t$ --- and are the mechanism through which
the quoter earns its inception spread. Limit orders rest inside the
book at a random offset from the mid,
\[
  p^{\mathrm{LO,\,buy}} \;=\; m_t\,\Bigl(1 - \tfrac{\varepsilon}{10^{4}}\Bigr),
  \qquad
  p^{\mathrm{LO,\,sell}} \;=\; m_t\,\Bigl(1 + \tfrac{\varepsilon}{10^{4}}\Bigr),
  \qquad
  \varepsilon \sim \mathrm{Exp}(k).
\]
Using the same $k$ in both the arrival intensity and the
limit-placement distribution is a self-consistency choice rather
than a fit: clients who are dissuaded from crossing at rate $k$ are
assumed to post passively at the same rate $k$. This keeps the
flow symmetric under a shift of the quoted distance --- widening our
spread by $\Delta\delta$ reduces arrivals by $e^{-k\Delta\delta}$
and simultaneously pushes the average passive limit deeper by the
same factor, with no second free parameter to tune.

\subsection{Heavy-tailed order sizes}
\label{sec:client-flow-size}

Client notional sizes are drawn from a truncated Pareto on the
interval $[x_{\min}, x_{\max}] = [10^{3}, 10^{5}]\,\mathrm{EUR}$
with tail index $\alpha = 1.5$, in line with the power-law volume
distribution of \citet{gabaix2003theory}. Truncation is not only
empirical realism --- the exchange rejects orders outside its
permitted notional band --- but also a numerical safeguard: an
untruncated Pareto with $\alpha = 1.5$ has infinite variance, and
Monte-Carlo P\&L moments would be dominated by a handful of
outliers instead of by the structural behaviour of the quoter. The
truncated inverse-CDF sampler is
\[
  X \;=\; x_{\min}\,\Bigl(1 - U\bigl(1 - (x_{\min}/x_{\max})^{\alpha}\bigr)\Bigr)^{-1/\alpha},
  \qquad U \sim \mathcal{U}(0,1),
\]
which respects the bounds exactly without rejection sampling. The
resulting distribution concentrates mass near $x_{\min}$ but
assigns a non-trivial weight to the largest permitted ticket, so a
typical session mixes many small fills with a few that can move
the inventory by several percent of capital in a single trade.

\begin{figure}[H]
  \centering
  \includegraphics[width=0.75\textwidth]{client-flow-pareto.png}
  \caption{Truncated Pareto size distribution on
           $[10^{3}, 10^{5}]$~EUR with $\alpha = 1.5$: mass
           concentrates near the lower bound while the heavy tail
           keeps a non-trivial weight at the hard cap $x_{\max}$.}
  \label{fig:client-flow-pareto}
\end{figure}

\subsection{Consumption by the order book}
\label{sec:client-flow-consumption}

Each generated order is tagged with $\mathrm{origin} = \texttt{client}$
and pushed into the same matching engine that ingests our own
quotes (Section~\ref{sec:orderbook}). A client market buy walks the
ask ladder level by level under price--time priority, partially
filling our resting quotes until its size is exhausted; a client
limit order either crosses and generates an immediate fill at the
touch we posted or rests in the book, where it remains eligible to
be matched against a future quote. The matching rule restricts
trades to client-versus-market-maker pairs, so client-to-client
crossings never occur and the realised fill price is always our
passive quote. This closed loop --- client intensity shaped by our
quote, fills generated only through our ladder --- is precisely the
structure assumed by the Avellaneda--Stoikov and \citet{gueant2013dealing}
half-spread derivations, and it makes the intensity baseline $A$
and decay $k$ the two knobs that jointly control fill rate and
adverse-selection severity alongside our posted half-spread $h_t$.
